\documentclass[12pt]{article}

% --- Packages ---
\usepackage[utf8]{inputenc}
\usepackage[T1]{fontenc}
\usepackage[french]{babel}
\usepackage{amsmath, amssymb}
\usepackage{geometry}
\usepackage{tikz}
\usepackage[most]{tcolorbox}
\usepackage{lmodern}
\usepackage{pgfplots}
\pgfplotsset{compat=1.18}
\usepackage{mdframed}
\usepackage{bm}
\usepackage{xcolor}
\usepackage{sectsty}
\usepackage[utf8]{inputenc}
\usepackage[T1]{fontenc}
\usepackage[dvipsnames]{xcolor}
\usepackage{tkz-tab}
\usepackage{multido}
% --- Page ---
\geometry{margin=2cm}

% Footer
\usepackage{fancyhdr}
\pagestyle{fancy}
\cfoot{\thepage}
\rfoot{\textit{G.M.}}
\renewcommand{\headrulewidth}{0pt}
\renewcommand{\footrulewidth}{0pt}

% Pour les normes : 
\newcommand*{\norme}[1]{\left\lVert\overrightarrow{#1}\right\rVert}

\usetikzlibrary{shapes.misc}

\tikzset{cross/.style={cross out, draw=black, minimum size=2*(#1-\pgflinewidth), inner sep=0pt, outer sep=0pt},
%default radius will be 1pt. 
cross/.default={4pt}}

% mettre dans le préambule, après \usepackage{xcolor}
\usepackage{sectsty} % gère facilement la couleur des titres
\setlength{\parindent}{0pt}

\sectionfont{\color{red}\Large\bfseries}     % pour \section*
% \subsectionfont{\color{red}\normalsize\bfseries} % pour \subsection*

\begin{document}

% ============================
%   Variables aléatoires
% ============================

\begin{center}
    {\LARGE \textbf{Variables aléatoires}}\\[4mm]
    {\large 1ère STMG}
\end{center}

\hrule
\vspace{0.4cm}

% ============================
%   PARTIE 1 : Variable aléatoire et loi de probabilité
% ============================

\section*{\textcolor{red}{Partie 1 : Variable aléatoire et loi de probabilité}}

\subsection*{\textcolor{red}{1.1 Variable aléatoire}}

\vspace{5pt}

\underline{Exemple :}
\vspace{10pt}
\\Soit l'expérience aléatoire : \textbf{« On lance un dé à six faces et on regarde le résultat. »}

L'ensemble de toutes les issues possibles $E = \{1, 2, 3, 4, 5, 6\}$ s'appelle l'univers des possibles.  

On considère le jeu suivant :  
\begin{itemize}
    \item Si le résultat est 5 ou 6, on gagne 2 €.
    \item Sinon, on perd 1 €.
\end{itemize}

On peut définir ainsi une variable aléatoire $X$ sur $E = \{1, 2, 3, 4, 5, 6\}$ qui donne le gain et qui 
peut prendre les valeurs $2$ ou $-1$.  

\[
X = 
\begin{cases}
2 & \text{si l'issue est \_\_\_\_\_\_\_\_\_\_\_}, \\[4mm]
-1 & \text{si l'issue est \_\_\_\_\_\_\_\_\_\_\_}.
\end{cases}
\]

\vspace{10pt}

\begin{mdframed}[
    linewidth=1pt,
    roundcorner=5pt,
    linecolor=ForestGreen
]
{\color{ForestGreen}\textbf{Définition :}}
Une \textbf{variable aléatoire} $X$ est une \underline{fonction} qui, à chaque issue d'une expérience 
aléatoire, associe un nombre réel.
\end{mdframed}

% \vspace{10pt}

\begin{mdframed}[
    linewidth=1pt,
    roundcorner=5pt,
    linecolor=RoyalBlue,
    innertopmargin=10pt,
    innerbottommargin=10pt
]
{\color{RoyalBlue}\textbf{\underline{Exercice d'application n°1 :} Calculer une probabilité à l’aide 
d’une variable aléatoire.}}
\vspace{5pt}
\\
Soit l'expérience aléatoire : « On tire une carte au hasard dans un jeu de 32 cartes. »  
\begin{itemize}
    \item Si cette carte est un cœur, on gagne 5 €.
    \item Si cette carte est un carreau, on gagne 2 €.
    \item Dans les autres cas, on perd 1 €.
\end{itemize}

Soit $X$ la variable aléatoire qui associe le gain du jeu.  
\vspace{5pt}
\\
Calculer : $P(X = 5)$, $P(X = -1)$ et $P(X \le 2)$.

\end{mdframed}

\vspace{20pt}

\subsection*{\textcolor{red}{1.2 Loi de probabilité}}

\vspace{5pt}

\begin{mdframed}[
    linewidth=1pt,
    roundcorner=5pt,
    linecolor=ForestGreen
]
{\color{ForestGreen}\textbf{Définition :}}
Définir la loi de probabilité d'une variable aléatoire $X$, c'est associer une 
probabilité à chaque valeur que peut prendre $X$.
\vspace{10pt}
\\La loi de probabilité d'une variable aléatoire est donnée, en général, 
dans un tableau à deux lignes :
\begin{itemize}
    \item la première donne les valeurs prises par $X$ (souvent rangées dans l'ordre 
    croissant)
    \item la seconde donne la probabilité que $X$ prenne chacune de ces valeurs
\end{itemize}
\vspace{10pt}
\[
\renewcommand{\arraystretch}{1.5}
\begin{array}{c|c|c|c|c}
x_i            & x_1 & x_2 & \cdots & x_n \\ \hline
P(X = x_i)     & p_1 & p_2 & \cdots & p_n
\end{array}
\]
\end{mdframed}

\underline{Remarque :} Les « $x_i$ » sont toutes les valeurs prises par $X$.

\vspace{10pt}
\begin{mdframed}[
    linewidth=1pt,
    roundcorner=5pt,
    linecolor=RoyalBlue,
    innertopmargin=10pt,
    innerbottommargin=10pt
]
{\color{RoyalBlue}\textbf{\underline{Exercice d'application n°2 :} Déterminer une loi de probabilité 
d’une variable aléatoire.}}
\vspace{5pt}
\\
On lance simultanément deux dés à 6 faces et on note les valeurs
obtenues.
\\Soit $X$ la variable aléatoire égale à la plus grande des deux valeurs.
\vspace{5pt}
\\Établir la loi de probabilité de $X$.

\end{mdframed}

\vspace{5pt}

% ============================
%   PARTIE 2 : Espérance
% ============================

\section*{\textcolor{red}{Partie 2 : Espérance}}

\begin{mdframed}[
    linewidth=1pt,
    roundcorner=5pt,
    linecolor=ForestGreen
]
{\color{ForestGreen}\textbf{Définition :}}
Soit une variable aléatoire $X$ prenant les valeurs $x_1, x_2, \dots, x_n$.\\
La loi de probabilité de $X$ associe à toute valeur $x_i$ la probabilité $p_i = P(X = x_i)$.\\
L'\textbf{espérance} de $X$ est :
\[
E(X) = p_1 x_1 + p_2 x_2 + \dots + p_n x_n
\]
\end{mdframed}

\vspace{5pt}

\begin{mdframed}[
    linewidth=1pt,
    roundcorner=5pt,
    linecolor=Red,
    innertopmargin=10pt,
    innerbottommargin=10pt
]
{\color{red}
\underline{Propriété :}
L'espérance d'une variable aléatoire est la valeur moyenne obtenue, par expérience,
au bout d'un \textbf{grand nombre de répétitions} de l'expérience.
}
\end{mdframed}

\vspace{5pt}

\begin{mdframed}[
    linewidth=1pt,
    roundcorner=5pt,
    linecolor=RoyalBlue,
    innertopmargin=10pt,
    innerbottommargin=10pt
]
{\color{RoyalBlue}\textbf{\underline{Exercice d'application n°3 :} Calculer d'une variable aléatoire.}}
\vspace{5pt}
\\
    On tire une carte dans un jeu de 32 cartes.
    \begin{itemize}
        \item Si on tire un cœur, on gagne 2 €.
        \item Si on tire un roi, on gagne 5 €.
        \item Si on tire une autre carte, on perd 1 €.
    \end{itemize}
    \vspace{5pt}
    $X$ est la variable aléatoire donnant le gain du jeu.
    \begin{enumerate}
        \item Calculer l'espérance de $X$.
        \item Donner une interprétation du résultat.
    \end{enumerate}
\end{mdframed}

\end{document}
